
Reinforcement learning from execution feedback for code generation requires non-zero reward variance within rollout groups in order for Group Relative Policy Optimization (GRPO) to produce informative policy gradients. Binary pass/fail rewards collapse to zero variance whenever a model cannot fully solve a problem, blocking all gradient flow. This paper analyzes the structural properties of ratio rewards --- the fraction of test cases passed --- relative to binary rewards under GRPO's group-relative advantage normalization. Under a Binomial model, the expected within-group variance of the ratio reward is E[Var(r\_ratio)] = q($1-q)/T$, while that of the binary reward is E[Var(r\_binary)] = $q^T$($1-q^{T})$, yielding a variance ratio of 5--$20\times$ in favor of ratio rewards for problems where 30--55\% of rollouts pass at least one test and T = 5. This analysis motivates a prescreening protocol that filters training problems to the partial-tractability window S\_term $\in$ [0.3, 0.55] before any gradient updates. A six-module prescreening pipeline was implemented and validated at 15/15 tasks and 67/67 integration tests. Applied to 300 APPS introductory problems with Qwen2.5-Coder-7B-Instruct (k=8 rollouts, 2,400 total solution attempts), prescreening yielded 0 qualifying problems. All 300 problems produced S\_term = 0.0: the base model passed zero test cases across all attempts. This result establishes that supervised fine-tuning initialization is a prerequisite for GRPO to produce non-degenerate training signal on APPS introductory problems with this model, a finding consistent with prior reports in the literature. All mechanism hypotheses (h-m1 through h-m4) remain unexecuted pending availability of an SFT checkpoint. The analytical variance advantage is mathematically established; its empirical consequences remain to be measured.

